% GENERATED FILE. Edit canonical lessons, then run npm run generate:books.
% canonical-source-hash: fnv1a64:2b44389244b7950c
% canonical-lessons: ML-C321-pulliyulla, ML-C321-sugandhamulla, ML-C321-narunna, ML-C321-bahalamulla, ML-C321-satyasandhamaya

\chapter{Spotted, Fragrant, Smelly, Noisy, Honest}
\label{ch:ml-spotted-fragrant-smelly-noisy-honest}
\hlchaptermodality{321}

\begin{chapteropening}
\textbf{By the end of this chapter:} \emph{I can say spotted, fragrant, smelly, noisy and honest.}

Complete the last lesson of chapter 321: I can say spotted, fragrant, smelly, noisy and honest.
\end{chapteropening}

\section[\texorpdfstring{\ml{പുള്ളിയുള്ള} (puḷḷiyuḷḷa)}{പുള്ളിയുള്ള (puḷḷiyuḷḷa)}]{\ml{പുള്ളിയുള്ള} (puḷḷiyuḷḷa) — spotted}

\label{lesson:ML-C321-pulliyulla}



\begin{quote}
Before the new one: say the Malayalam for terrible, then the Malayalam for awesome.
\end{quote}

\subsection*{You'll want to know: \ml{പുള്ളിയുള്ള}}
\textbf{\ml{പുള്ളിയുള്ള}} — \emph{puḷḷiyuḷḷa} — \textquotedblleft{}spotted\textquotedblright{}.

\begin{cousinweb}
\textbf{\ml{പുള്ളിയുള്ള}} (\emph{puḷḷiyuḷḷa}) is native Dravidian, \textquotedblleft{}having spots\textquotedblright{}: \emph{puḷḷi}, \textquotedblleft{}dot, spot\textquotedblright{} (DEDR 4327, with Tamil \emph{puḷḷi}), and \emph{-uḷḷa}, \textquotedblleft{}that has\textquotedblright{} (DEDR 697).
\end{cousinweb}

\begin{grammarlens}[title={What you've built}]
One more word for this chapter.
\end{grammarlens}

\subsection*{Your turn}
\begin{itemize}
\raggedright
  \item \emph{Say it:} \emph{puḷḷiyuḷḷa}
  \item \emph{Say it:} \emph{puḷḷiyuḷḷa}, once more
  \item \emph{Say it:} say \emph{puḷḷiyuḷḷa}
  \item \emph{Recall:} say the Malayalam for terrible, then the Malayalam for awesome, then say \emph{puḷḷiyuḷḷa} again
\end{itemize}

\begin{tcolorbox}[breakable,colback=teal!4,colframe=teal!35!black,title={Before you move on}]
What does \ml{പുള്ളിയുള്ള} mean? (\textquotedblleft{}Spotted\textquotedblright{}.) Say it once more.
\end{tcolorbox}
\section[\texorpdfstring{\ml{സുഗന്ധമുള്ള} (sugandhamuḷḷa)}{സുഗന്ധമുള്ള (sugandhamuḷḷa)}]{\ml{സുഗന്ധമുള്ള} (sugandhamuḷḷa) — fragrant}

\label{lesson:ML-C321-sugandhamulla}



\begin{quote}
Before the new one: say the Malayalam for awesome, then the Malayalam for spotted.
\end{quote}

\subsection*{You'll want to know: \ml{സുഗന്ധമുള്ള}}
\textbf{\ml{സുഗന്ധമുള്ള}} — \emph{sugandhamuḷḷa} — \textquotedblleft{}fragrant\textquotedblright{}.

\begin{cousinweb}
\textbf{\ml{സുഗന്ധമുള്ള}} (\emph{sugandhamuḷḷa}) joins Sanskrit \emph{sugandha}, \textquotedblleft{}fragrant, perfume\textquotedblright{} (CDIAL 13454): \emph{su}, \textquotedblleft{}good\textquotedblright{}, and \emph{gandha}, \textquotedblleft{}smell\textquotedblright{}, to native \emph{-uḷḷa}, \textquotedblleft{}that has\textquotedblright{} (DEDR 697).
\end{cousinweb}

\begin{grammarlens}[title={What you've built}]
One more word for this chapter.
\end{grammarlens}

\subsection*{Your turn}
\begin{itemize}
\raggedright
  \item \emph{Say it:} \emph{sugandhamuḷḷa}
  \item \emph{Say it:} \emph{sugandhamuḷḷa}, once more
  \item \emph{Say it:} say \emph{sugandhamuḷḷa}
  \item \emph{Recall:} say the Malayalam for awesome, then the Malayalam for spotted, then say \emph{sugandhamuḷḷa} again
\end{itemize}

\begin{tcolorbox}[breakable,colback=teal!4,colframe=teal!35!black,title={Before you move on}]
What does \ml{സുഗന്ധമുള്ള} mean? (\textquotedblleft{}Fragrant\textquotedblright{}.) Say it once more.
\end{tcolorbox}
\section[\texorpdfstring{\ml{നാറുന്ന} (nāṟunna)}{നാറുന്ന (nāṟunna)}]{\ml{നാറുന്ന} (nāṟunna) — smelly, stinking}

\label{lesson:ML-C321-narunna}



\begin{quote}
Before the new one: say the Malayalam for spotted, then the Malayalam for fragrant.
\end{quote}

\subsection*{You'll want to know: \ml{നാറുന്ന}}
\textbf{\ml{നാറുന്ന}} — \emph{nāṟunna} — \textquotedblleft{}smelly, stinking\textquotedblright{}.

\begin{cousinweb}
\textbf{\ml{നാറുന്ന}} (\emph{nāṟunna}) is native Dravidian, a form of the verb \textquotedblleft{}to give off a smell, stink\textquotedblright{}; DEDR 2918 lists the root with Tamil \emph{ñāṟu} and Kannada \emph{nāṟu}.
\end{cousinweb}

\begin{grammarlens}[title={What you've built}]
One more word for this chapter.
\end{grammarlens}

\subsection*{Your turn}
\begin{itemize}
\raggedright
  \item \emph{Say it:} \emph{nāṟunna}
  \item \emph{Say it:} \emph{nāṟunna}, once more
  \item \emph{Say it:} say \emph{nāṟunna}
  \item \emph{Recall:} say the Malayalam for spotted, then the Malayalam for fragrant, then say \emph{nāṟunna} again
\end{itemize}

\begin{tcolorbox}[breakable,colback=teal!4,colframe=teal!35!black,title={Before you move on}]
What does \ml{നാറുന്ന} mean? (\textquotedblleft{}Smelly, stinking\textquotedblright{}.) Say it once more.
\end{tcolorbox}
\section[\texorpdfstring{\ml{ബഹളമുള്ള} (bahaḷamuḷḷa)}{ബഹളമുള്ള (bahaḷamuḷḷa)}]{\ml{ബഹളമുള്ള} (bahaḷamuḷḷa) — noisy}

\label{lesson:ML-C321-bahalamulla}



\begin{quote}
Before the new one: say the Malayalam for fragrant, then the Malayalam for smelly.
\end{quote}

\subsection*{You'll want to know: \ml{ബഹളമുള്ള}}
\textbf{\ml{ബഹളമുള്ള}} — \emph{bahaḷamuḷḷa} — \textquotedblleft{}noisy\textquotedblright{}.

\begin{cousinweb}
\textbf{\ml{ബഹളമുള്ള}} (\emph{bahaḷamuḷḷa}) is \textquotedblleft{}having noise\textquotedblright{}: \emph{bahaḷaṁ}, \textquotedblleft{}uproar\textquotedblright{}, a word of uncertain origin, with native \emph{-uḷḷa}, \textquotedblleft{}that has\textquotedblright{} (DEDR 697).
\end{cousinweb}

\begin{grammarlens}[title={What you've built}]
One more word for this chapter.
\end{grammarlens}

\subsection*{Your turn}
\begin{itemize}
\raggedright
  \item \emph{Say it:} \emph{bahaḷamuḷḷa}
  \item \emph{Say it:} \emph{bahaḷamuḷḷa}, once more
  \item \emph{Say it:} say \emph{bahaḷamuḷḷa}
  \item \emph{Recall:} say the Malayalam for fragrant, then the Malayalam for smelly, then say \emph{bahaḷamuḷḷa} again
\end{itemize}

\begin{tcolorbox}[breakable,colback=teal!4,colframe=teal!35!black,title={Before you move on}]
What does \ml{ബഹളമുള്ള} mean? (\textquotedblleft{}Noisy\textquotedblright{}.) Say it once more.
\end{tcolorbox}
\section[\texorpdfstring{\ml{സത്യസന്ധമായ} (satyasandhamāya)}{സത്യസന്ധമായ (satyasandhamāya)}]{\ml{സത്യസന്ധമായ} (satyasandhamāya) — honest}

\label{lesson:ML-C321-satyasandhamaya}



\begin{quote}
Before the new one: say the Malayalam for smelly, then the Malayalam for noisy.
\end{quote}

\subsection*{You'll want to know: \ml{സത്യസന്ധമായ}}
\textbf{\ml{സത്യസന്ധമായ}} — \emph{satyasandhamāya} — \textquotedblleft{}honest\textquotedblright{}.

\begin{cousinweb}
\textbf{\ml{സത്യസന്ധമായ}} (\emph{satyasandhamāya}) is Sanskrit \emph{satyasandha}, \textquotedblleft{}true to one's promise\textquotedblright{}: \emph{satya}, \textquotedblleft{}truth\textquotedblright{} (CDIAL 13112), and \emph{sandhā}, \textquotedblleft{}a promise\textquotedblright{}, with native \emph{āya}, \textquotedblleft{}having become\textquotedblright{} (DEDR 333). DEDR 333 compares Sanskrit \emph{ām}.
\end{cousinweb}

\begin{grammarlens}[title={What you've built}]
One more word for this chapter.
\end{grammarlens}

\subsection*{Your turn}
\begin{itemize}
\raggedright
  \item \emph{Say it:} \emph{satyasandhamāya}
  \item \emph{Say it:} \emph{satyasandhamāya}, once more
  \item \emph{Say it:} say \emph{satyasandhamāya}
  \item \emph{Recall:} say the Malayalam for smelly, then the Malayalam for noisy, then say \emph{satyasandhamāya} again
\end{itemize}

\begin{tcolorbox}[breakable,colback=teal!4,colframe=teal!35!black,title={Before you move on}]
What does \ml{സത്യസന്ധമായ} mean? (\textquotedblleft{}Honest\textquotedblright{}.) Say it once more.
\end{tcolorbox}
